% THEORY.tex
% Proof ledger for the project.  This file records mathematical truth and proof
% status; it does NOT determine how many theorems appear in the main paper.
%
% STATUS LEGEND
%   PROVED      = proof is complete under the stated assumptions.
%   PROOF TASK  = theorem route is fixed but details/constants remain to write.
%   BLOCKER     = validity for the intended paper depends on resolving the item.
%
% CURRENT STATUS
%   T1 Exact gap transport:                     PROVED
%   T2 Screenable omitted-factor bound:          PROVED (abstract form)
%   T2 reflection-symmetry construction:          PROVED
%   T2 nonlinear-PDE construction:                BLOCKER
%   T3 Full-action certificate:                  PROVED
%   T4 Decision/conditioning separation:         PROVED under abstract assumptions;
%                                                 concrete mapping is a PROOF TASK
%   Rigorous inference instantiation:            BLOCKER for an end-to-end certified demo

\section{Theory Ledger}
\label{sec:theory-ledger}

\paragraph{Notation.}
All notation follows \texttt{PROBLEM\_FORMULATION.tex}.  In particular,
$\pi_S$ is the active conditioned law, $\pi_C$ the fully conditioned law,
$F_{x,\widehat x}=u_x-u_{\widehat x}$ the acquisition-gap observable, and
$G_Q(x,\widehat x)=\mathbb E_Q[F_{x,\widehat x}]$.  For a factor-sufficient
Gaussian vector $Y$, $\bar F_{x,\widehat x}
=\mathbb E_{P_{0,t}}[F_{x,\widehat x}\mid Y]$ denotes the Rao--Blackwellized
gap observable.

% -----------------------------------------------------------------------------
\subsection{T1: Exact acquisition-gap transport}

\paragraph{Purpose.}
Answer the BO question: \emph{how can omitted conditioning change the ordering
between the proposed experiment and a challenger?}

\paragraph{Status.}
\textbf{PROVED.}  This is a standard exponential-tilting derivative specialized
to a BO acquisition gap.  The identity is central; the calculus is not claimed
as novel.

\paragraph{Assumptions.}
Fix an active set $S$, challenger $x$, and proposed action $\widehat x$.  Let
\[
    H_S(f)=\sum_{j\notin S}e_j(f),
\]
and define
\[
    \pi_{S,s}(df)
    = Z_s^{-1}\exp[-E_S(f)-sH_S(f)]P_{0,t}(df),
    \qquad s\in[0,1].
\]
Assume:
\begin{enumerate}
    \item $0<Z_s<\infty$ for all $s\in[0,1]$;
    \item $F_{x,\widehat x}$, $H_S$, and
          $F_{x,\widehat x}H_S$ are integrable under $\pi_{S,s}$;
    \item differentiation under the integral is valid on $[0,1]$ (for example,
          by an integrable dominating function in a neighborhood of the path).
\end{enumerate}
No differentiability of the BO utility is required for this result.

\paragraph{Theorem T1 (exact acquisition-gap transport).}
Let
\[
    G_s(x,\widehat x)
    =\mathbb E_{\pi_{S,s}}[F_{x,\widehat x}].
\]
Then $G_s$ is absolutely continuous and
\begin{equation}
    \frac{d}{ds}G_s(x,\widehat x)
    =
    -\operatorname{Cov}_{\pi_{S,s}}
       \!\left(F_{x,\widehat x},H_S\right).
    \label{eq:T1-derivative}
\end{equation}
Consequently,
\begin{equation}
    G_C(x,\widehat x)-G_S(x,\widehat x)
    =
    -\int_0^1
      \operatorname{Cov}_{\pi_{S,s}}
      \!\left(F_{x,\widehat x},H_S\right)\,ds.
    \label{eq:T1-integrated}
\end{equation}
If $H_S=\sum_{j\notin S}e_j$, then
\begin{equation}
    G_C-G_S
    =
    -\sum_{j\notin S}\int_0^1
      \operatorname{Cov}_{\pi_{S,s}}
      \!\left(F_{x,\widehat x},e_j\right)\,ds.
    \label{eq:T1-factorized}
\end{equation}

\begin{proof}
Write
\[
    N_s
    = \int F_{x,\widehat x}(f)
      \exp[-E_S(f)-sH_S(f)]\,P_{0,t}(df),
\]
so that $G_s=N_s/Z_s$.  Differentiation under the integral gives
\begin{align*}
    N_s'
    &= -Z_s\,\mathbb E_{\pi_{S,s}}[F_{x,\widehat x}H_S],\\
    Z_s'
    &= -Z_s\,\mathbb E_{\pi_{S,s}}[H_S].
\end{align*}
The quotient rule therefore yields
\begin{align*}
    G_s'
    &= \frac{N_s'Z_s-N_sZ_s'}{Z_s^2}\\
    &= -\mathbb E_{\pi_{S,s}}[F_{x,\widehat x}H_S]
       +\mathbb E_{\pi_{S,s}}[F_{x,\widehat x}]
        \mathbb E_{\pi_{S,s}}[H_S]\\
    &= -\operatorname{Cov}_{\pi_{S,s}}
       (F_{x,\widehat x},H_S).
\end{align*}
Integrating from $s=0$ to $s=1$ proves
Equation~\eqref{eq:T1-integrated}.  The factorized form follows from covariance
linearity and finite summation.
\end{proof}

\paragraph{Main-paper use.}
State the theorem and include the short proof or derivation.  This is the
conceptual bridge from rich conditioning to BO decisions.

% -----------------------------------------------------------------------------
\subsection{T2: Screenable omitted-factor influence}

\paragraph{Purpose.}
Turn T1 from an exact but expensive identity into a computable upper bound that
does not require inference under every omitted factor.

\paragraph{Status.}
\textbf{PROVED in abstract form.}  Proposition~T2-SYM below proves the concrete
reflection-symmetry construction.  The nonlinear-PDE construction remains a
separate \textbf{BLOCKER}.

\paragraph{Screenability assumption S.}
Assume that, for every active set encountered by the algorithm and every
$s\in[0,1]$, the relevant full-function observable $g$ and factor-measurable
observable $h$ satisfy
\begin{equation}
    |\operatorname{Cov}_{\pi_{S,s}}(g,h)|
    \le L(\bar g)^\top C_S L(h),
    \label{eq:S-cov}
\end{equation}
where $L(\bar g)\in\mathbb R_+^B$ is a computable block-sensitivity vector and
$C_S$ is entrywise nonnegative and available without sampling from $\pi_C$.

\paragraph{Theorem T2 (omitted-factor influence bound).}
Under the assumptions of T1 and Assumption S,
\begin{equation}
    |G_C(x,\widehat x)-G_S(x,\widehat x)|
    \le
    B_{\mathrm{struct}}(x,\widehat x),
    \label{eq:T2-bound}
\end{equation}
with
\begin{equation}
    B_{\mathrm{struct}}(x,\widehat x)
    =
    L(\bar F_{x,\widehat x})^\top C_S
    \sum_{j\notin S}L(e_j).
    \label{eq:T2-Bstruct}
\end{equation}
If $C_S=A_S^{-1}$, define
\[
    h_U=\sum_{j\notin S}L(e_j)
\]
and solve
\[
    A_S w=h_U.
\]
Then
\[
    B_{\mathrm{struct}}(x,\widehat x)
    =L(\bar F_{x,\widehat x})^\top w
\]
for every challenger $x$.

\begin{proof}
By T1 and the triangle inequality,
\begin{align*}
    |G_C-G_S|
    &\le
    \sum_{j\notin S}\int_0^1
    \left|
    \operatorname{Cov}_{\pi_{S,s}}
    (F_{x,\widehat x},e_j)
    \right|ds.
\end{align*}
Apply Assumption S to each term:
\begin{align*}
    |G_C-G_S|
    &\le
    \sum_{j\notin S}\int_0^1
    L(\bar F_{x,\widehat x})^\top C_S L(e_j)\,ds\\
    &=
    L(\bar F_{x,\widehat x})^\top C_S
    \sum_{j\notin S}L(e_j),
\end{align*}
which is Equation~\eqref{eq:T2-Bstruct}.  The linear-solve form is immediate
when $C_S=A_S^{-1}$.
\end{proof}

\paragraph{Published covariance theorem.}
The concrete reflection-symmetry construction invokes Theorem~2.3
(``Covariance estimate, Otto \& Menz'') of Georg Menz,
\emph{A Brascamp--Lieb type covariance estimate}, Electron. J. Probab. 19
(2014), no.~78, 1--15, doi:10.1214/EJP.v19-2997.  For a measure
$\mu(dy)\propto\exp[-V(y)]dy$ on a product of Euclidean blocks, Menz assumes
that $V$ is convex at infinity, that every one-block conditional law has a
uniform Poincar\'e constant $\rho_i>0$, that
$\|\nabla_i\nabla_jV\|_{\mathrm{op}}\le\kappa_{ij}$ uniformly for $i\ne j$,
and that
\[
    A_{ii}=\rho_i,
    \qquad
    A_{ij}=-\kappa_{ij}\quad(i\ne j)
\]
is positive definite.  The theorem gives
\begin{equation}
 |\operatorname{Cov}_{\mu}(g,h)|
 \le
 \sum_{i,j}(A^{-1})_{ij}
 \left(\mathbb E_{\mu}\|\nabla_i g\|_2^2\right)^{1/2}
 \left(\mathbb E_{\mu}\|\nabla_j h\|_2^2\right)^{1/2}.
 \label{eq:menz-covariance}
\end{equation}
The proof also establishes that $A^{-1}$ is entrywise nonnegative.

\paragraph{Proposition T2-SYM (uniform reflection-symmetry influence).}
Fix one BO iteration.  Let
$Y=(Y_1,\ldots,Y_N)\in(\mathbb R^2)^N$ have reference law
$\mathcal N(m,K)$ with precision $Q=K^{-1}\succ0$, where
\[
    Y_j=\begin{bmatrix}f(-r_j)\\f(r_j)\end{bmatrix},
    \qquad
    b=\begin{bmatrix}-1\\1\end{bmatrix}.
\]
For $\gamma_j,\tau_j>0$, define the block-local factors
\[
    e_j(Y_j)
    =\gamma_j\log\cosh\!\left(\frac{b^\top Y_j}{\tau_j}\right).
\]
For $S\subseteq[N]$ and $s\in[0,1]$, let $\mu_{S,s}$ have potential
\begin{equation}
 V_{S,s}(y)
 =\frac12(y-m)^\top Q(y-m)
 +\sum_{j=1}^N
 \left(\mathbf 1\{j\in S\}+s\mathbf 1\{j\notin S\}\right)e_j(y_j).
 \label{eq:symmetry-potential}
\end{equation}
Write $Q_{ij}$ for its two-by-two block entries and define
\begin{equation}
 \rho_i=\lambda_{\min}(Q_{ii}),
 \qquad
 \kappa_{ij}=\|Q_{ij}\|_{\mathrm{op}},
 \qquad
 A_{ii}=\rho_i,
 \quad
 A_{ij}=-\kappa_{ij}\ (i\ne j).
 \label{eq:symmetry-comparison}
\end{equation}
If $A\succ0$, then the single conservative choice
\[
    C_S\equiv C=A^{-1}
\]
is valid simultaneously for every active set $S$ and every $s\in[0,1]$.
Moreover,
\begin{equation}
    L_i(e_j)=\mathbf 1\{i=j\}\frac{\sqrt2\,\gamma_j}{\tau_j}.
    \label{eq:symmetry-factor-sensitivity}
\end{equation}

Suppose the factors are measurable with respect to $Y$ and the reference GP
conditional law is
\[
 f(x)\mid Y=y
 \sim
 \mathcal N\!\left(m(x)+a_x^\top(y-m),\sigma_x^2\right),
 \qquad
 a_x=K^{-1}k_{Yx}.
\]
For EI, let
\[
 \bar u_x(y)=\mathbb E[(f(x)-y_t^\star)_+\mid Y=y],
 \qquad
 \bar F_{x,\widehat x}=\bar u_x-\bar u_{\widehat x},
\]
and partition $a_x$ into the same two-coordinate blocks.  A valid block
sensitivity is
\begin{equation}
 d_i^{\mathrm{EI}}(x,\widehat x)
 =\max\!\left\{
 \|a_{x,i}\|_2,
 \|a_{\widehat x,i}\|_2,
 \|a_{x,i}-a_{\widehat x,i}\|_2
 \right\}.
 \label{eq:ei-block-footprint}
\end{equation}
Consequently, for omitted factors $U=S^c$, with
\[
 (h_U)_i
 =\mathbf 1\{i\notin S\}\frac{\sqrt2\,\gamma_i}{\tau_i},
\]
the concrete structural bound is
\begin{equation}
 |G_C(x,\widehat x)-G_S(x,\widehat x)|
 \le
 \bigl(d^{\mathrm{EI}}(x,\widehat x)\bigr)^\top A^{-1}h_U.
 \label{eq:symmetry-ei-bound}
\end{equation}

\begin{proof}
The symmetry-factor derivatives are
\[
 \nabla_j e_j
 =\frac{\gamma_j}{\tau_j}
 \tanh\!\left(\frac{b^\top y_j}{\tau_j}\right)b,
 \qquad
 \nabla_j^2e_j
 =\frac{\gamma_j}{\tau_j^2}
 \operatorname{sech}^2\!\left(\frac{b^\top y_j}{\tau_j}\right)bb^\top
 \succeq0.
\]
This proves Equation~\eqref{eq:symmetry-factor-sensitivity}.  Conditional on
$y_{-i}$, the $i$th one-block potential has Hessian
\[
 Q_{ii}
 +\left(\mathbf 1\{i\in S\}+s\mathbf 1\{i\notin S\}\right)
   \nabla_i^2e_i
 \succeq Q_{ii}.
\]
The Bakry--\'Emery criterion therefore gives a conditional Poincar\'e constant
at least $\rho_i=\lambda_{\min}(Q_{ii})$, uniformly in the conditioning value,
$S$, and $s$.  Factor locality gives
$\nabla_i\nabla_jV_{S,s}=Q_{ij}$ for $i\ne j$.  Finally,
$\nabla^2V_{S,s}\succeq Q\succ0$, so the potential is globally strongly
convex and in particular satisfies Menz's convex-at-infinity assumption.
Menz's Theorem~2.3 now gives Equation~\eqref{eq:menz-covariance}; bounding its
block $L^2$ gradient norms by essential suprema yields
\[
 |\operatorname{Cov}_{\mu_{S,s}}(g,h)|
 \le L(g)^\top A^{-1}L(h)
\]
uniformly over $S$ and $s$.

Because all reweighting factors depend only on $Y$,
Equation~\eqref{eq:rao-blackwell-covariance} transfers the required covariance
to the finite-dimensional marginal.  If $\sigma_x>0$, conditional EI has
\[
 \bar u_x(y)
 = (\mu_x(y)-y_t^\star)\Phi(z_x)+\sigma_x\phi(z_x),
 \qquad
 z_x=\frac{\mu_x(y)-y_t^\star}{\sigma_x},
\]
and $\nabla_i\bar u_x=\Phi(z_x)a_{x,i}$.  Hence
$\nabla_i\bar F=p a_{x,i}-q a_{\widehat x,i}$ for some $p,q\in[0,1]$.
This vector is a convex combination of
$0$, $a_{x,i}$, $-a_{\widehat x,i}$, and
$a_{x,i}-a_{\widehat x,i}$, so convexity of the norm proves
Equation~\eqref{eq:ei-block-footprint}.

When $\sigma_x=0$, conditional EI is the positive-part function.  Its weak
derivative has a multiplier in $[0,1]$ almost everywhere.  Equivalently,
approximate the positive part by smooth 1-Lipschitz functions whose derivatives
remain in $[0,1]$, apply Menz's smooth result, and pass to the limit using the
Gaussian tails.  Thus the same footprint holds at the nonsmooth limit.
Substitution into Theorem~T2 gives Equation~\eqref{eq:symmetry-ei-bound}.
\end{proof}

\paragraph{Analytic positivity for the archived OU model.}
For the equally spaced OU kernel with
$q=\exp(-\Delta x/\ell)\in(0,1)$, the spatial precision is the AR(1)
tridiagonal matrix with endpoint diagonal $1/(1-q^2)$, interior diagonal
$(1+q^2)/(1-q^2)$, and adjacent magnitude $q/(1-q^2)$.  In inner-to-outer
reflection-block order, $A$ is tridiagonal.  For $N\ge2$, its minimum strict
row-diagonal-dominance margin is
\[
 \frac{1-q}{1+q}>0.
\]
Gershgorin's theorem therefore proves $A\succ0$ analytically for every such
$N$; a numerical eigenvalue is only a regression diagnostic.

\paragraph{Relationship to the archived log-acquisition prototype.}
The archived symmetry notebook does not use EI.  It uses exponential utility
$u_x(f)=\exp(\beta f(x))$ and bounds a difference of log acquisitions.  Its
footprint
\[
 d_i=\beta\|[a_x-a_{\widehat x}]_i\|_2
\]
is justified by an additional interpolation between the two action-dependent
linear tilts.  Specifically, use the potential
\[
 V_{S,s,r}^{x,\widehat x}(y)
 =V_{S,s}(y)
 -\beta\{(1-r)a_{\widehat x}+r a_x\}^\top(y-m),
 \qquad r\in[0,1].
\]
The added term is linear, so the same $A$ is valid uniformly in $r$ as well as
$S,s$.  Differentiation in $r$, followed by Menz's covariance bound and
integration over both unit intervals, gives the notebook's structural term
$h_U^\top A^{-1}d$ with no extra multiplicative factor.  This validates that
historical structural calculation; it does not relabel the prototype as EI or
supply its separate Monte Carlo and action-grid errors.

\paragraph{Reflection-symmetry status and remaining blocker.}
The proposition resolves the analytic T2-B construction for the stated
reflection-symmetry family whenever $A\succ0$, including the archived OU
specialization above.  It does not prove the nonlinear-PDE construction.  The
rigorous inference and action-optimization terms required for an end-to-end T3
certificate also remain separate.

\paragraph{Main-paper use.}
State Proposition~T2-SYM and Equation~\eqref{eq:symmetry-ei-bound} concisely.
Place the Menz assumption check, OU precision calculation, weak-derivative
argument, and archived-prototype distinction in the supplement.

% -----------------------------------------------------------------------------
\subsection{T3: Full-target BO action certificate}

\paragraph{Purpose.}
Answer the stopping question: \emph{when is the current experiment guaranteed to
be within $\epsilon$ of the experiment selected by fully conditioned BO?}

\paragraph{Status.}
\textbf{PROVED.}  The theorem is elementary once valid structural, inference,
and optimization error bounds are supplied.  The experimental challenge is
constructing those bounds rigorously in at least one complete implementation.

\paragraph{Inference interface.}
Let $\widehat G_S(x,\widehat x)$ satisfy, on an event $\mathcal E$ with
$\Pr(\mathcal E)\ge1-\delta$,
\begin{equation}
    |\widehat G_S(x,\widehat x)-G_S(x,\widehat x)|
    \le B_{\mathrm{infer}}(x,\widehat x)
    \label{eq:T3-infer}
\end{equation}
simultaneously over the challengers used by the certification procedure.

Define
\begin{equation}
    \Psi_S(x;\widehat x)
    = \widehat G_S(x,\widehat x)
      +B_{\mathrm{infer}}(x,\widehat x)
      +B_{\mathrm{struct}}(x,\widehat x).
    \label{eq:T3-Psi}
\end{equation}
Suppose a challenger routine returns $x_c$ and $\eta_{\mathrm{opt}}\ge0$ such
that
\begin{equation}
    \sup_{x\in\mathcal X}\Psi_S(x;\widehat x)
    \le
    \Psi_S(x_c;\widehat x)+\eta_{\mathrm{opt}}.
    \label{eq:T3-opt}
\end{equation}

\paragraph{Theorem T3 (full-conditioned action certificate).}
Under T2, Equations~\eqref{eq:T3-infer}--\eqref{eq:T3-opt} imply that, on
$\mathcal E$,
\begin{equation}
    R_C(\widehat x)
    \le
    \Psi_S(x_c;\widehat x)+\eta_{\mathrm{opt}}.
    \label{eq:T3-certificate}
\end{equation}
Hence, if the right-hand side is at most $\epsilon$, then
\[
    \alpha_C(x_C^\star)-\alpha_C(\widehat x)\le\epsilon.
\]

\begin{proof}
For every challenger $x$, T2 gives
\[
    G_C(x,\widehat x)
    \le G_S(x,\widehat x)+B_{\mathrm{struct}}(x,\widehat x).
\]
On $\mathcal E$, Equation~\eqref{eq:T3-infer} gives
\[
    G_S(x,\widehat x)
    \le \widehat G_S(x,\widehat x)+B_{\mathrm{infer}}(x,\widehat x).
\]
Therefore
\[
    G_C(x,\widehat x)\le\Psi_S(x;\widehat x).
\]
Taking the supremum over $x$ and applying Equation~\eqref{eq:T3-opt},
\begin{align*}
    R_C(\widehat x)
    &=\sup_x G_C(x,\widehat x)\\
    &\le\sup_x\Psi_S(x;\widehat x)\\
    &\le\Psi_S(x_c;\widehat x)+\eta_{\mathrm{opt}}.
\end{align*}
\end{proof}

\paragraph{Useful corollary: leader optimization need not be certified.}
T3 holds for any returned $\widehat x$.  The leader optimizer can therefore be
approximate or local.  A poor leader simply leaves a positive optimistic
challenger and prevents certification.  The globally certified optimization
requirement belongs to the challenger envelope, not to the initial leader.

\paragraph{Main-paper use.}
State and prove this result.  It is the formal reason the method can return an
action without all pairwise action comparisons.

% -----------------------------------------------------------------------------
\subsection{T4: When decision-relevant conditioning stays small}

\paragraph{Purpose.}
Formalize the empirical phenomenon that the conditioning required to settle the
next BO action can remain local even while the total factorized conditioning
problem grows.

\paragraph{Status.}
\textbf{PROVED under the abstract assumptions below.}  Mapping the theorem
cleanly onto the concrete PDE/symmetry models, including inverse-decay constants,
is a \textbf{PROOF TASK}.  The lower-bound witness concerns subsets of the
\emph{original factors}; do not overstate it as an impossibility result against
all posterior approximations.

\paragraph{Definition.}
For original-factor subsets, define
\begin{equation}
    M_{\mathrm{dec}}(\epsilon)
    =
    \min\left\{|S|:
    \exists\,x_S\ \text{with}\ R_C(x_S)\le\epsilon
    \ \text{certifiable using }S\right\}.
    \label{eq:Mdec}
\end{equation}
For a posterior-fidelity comparison, one possible benchmark is
\begin{equation}
    M_{\mathrm{post}}(\delta)
    =
    \min\left\{|S|:
    D_{\mathrm{KL}}(\pi_C\|\pi_S)\le\delta\right\}.
    \label{eq:Mpost}
\end{equation}
This second definition is only one posterior-oriented reference criterion.

\paragraph{Assumptions for localized decision influence.}
Consider a sequence of factor graphs with block set $V_n$ and graph distance
$d(\cdot,\cdot)$.  Let $D\subset V_n$ be a fixed-size decision region.
Assume:
\begin{enumerate}
    \item \textbf{Localized gap sensitivity:}
          $L_i(\bar F_{x,\widehat x})=0$ for $i\notin D$ and
          $\sup_{x,\widehat x}\|L(\bar F_{x,\widehat x})\|_1\le L_F$.
    \item \textbf{Bounded factor density/load:} each factor has uniformly
          bounded support diameter, and for every block $k$,
          \[
              \sum_{j:k\in\operatorname{supp}(e_j)}L_k(e_j)\le H_0.
          \]
    \item \textbf{Exponential influence decay:}
          \[
              (C_S)_{ik}\le c_C\exp[-\beta d(i,k)]
          \]
          for constants independent of the total domain size.
    \item \textbf{Polynomial graph growth:} the number of blocks at distance
          exactly $\ell$ from the fixed decision region is bounded by
          \[
              c_V(1+\ell)^{d_g-1}.
          \]
    \item \textbf{Bounded factor density in neighborhoods:} the number of
          factors whose support intersects a radius-$r$ neighborhood of $D$ is
          at most $c_M(1+r)^{d_g}$.
\end{enumerate}

\paragraph{Theorem T4A (localized conditioning upper bound).}
Activate every factor whose support intersects the radius-$r$ neighborhood of
$D$.  Then there is a constant $K$, independent of the total number of factors,
such that
\begin{equation}
    \sup_x B_{\mathrm{struct}}(x,\widehat x)
    \le
    K(1+r)^{d_g-1}e^{-\beta r}.
    \label{eq:T4-tail}
\end{equation}
Consequently, there exists
\[
    r_\epsilon=O\!\left(\log\frac{1}{\epsilon}\right)
\]
(up to the logarithmic polynomial correction in
Equation~\eqref{eq:T4-tail}) for which the structural error is at most
$\epsilon$, and
\begin{equation}
    M_{\mathrm{dec}}(\epsilon)
    =
    O\!\left(\log^{d_g}\frac{1}{\epsilon}\right),
    \label{eq:T4-Mdec}
\end{equation}
with constants independent of the total conditioning-domain size.

\begin{proof}
Because the acquisition-gap sensitivity is supported on $D$,
Equation~\eqref{eq:T2-Bstruct} and the factor-load assumption give
\begin{align*}
B_{\mathrm{struct}}(x,\widehat x)
&\le
\sum_{i\in D}L_i(\bar F_{x,\widehat x})
\sum_{k:d(k,D)>r}
(C_S)_{ik}
\sum_{j:k\in\operatorname{supp}(e_j)}L_k(e_j)\\
&\le
L_F c_C H_0
\sum_{\ell>r}c_V(1+\ell)^{d_g-1}e^{-\beta\ell}.
\end{align*}
The tail of an exponentially weighted polynomial sequence is bounded by
$K'(1+r)^{d_g-1}e^{-\beta r}$ for a constant $K'$ depending only on
$(\beta,d_g)$.  This proves Equation~\eqref{eq:T4-tail}.  Choosing $r$ so that
this bound is at most $\epsilon$ gives
$r_\epsilon=O(\log(1/\epsilon))$ up to the polynomial-log correction.  The
neighborhood factor-density assumption then gives
$M_{\mathrm{dec}}(\epsilon)\le c_M(1+r_\epsilon)^{d_g}$.
\end{proof}

\paragraph{Theorem T4B (posterior-fidelity lower-bound witness).}
Consider $N$ independent latent blocks.  Let the fully conditioned target
factorize as
\[
    \pi_C=\bigotimes_{j=1}^N q_j,
\]
and let omitting factor $j$ leave a reference marginal $\gamma_j$.  Suppose
\[
    D_{\mathrm{KL}}(q_j\|\gamma_j)\ge\kappa>0
\]
for all $j$.  For an original-factor subset $S$,
\[
    \pi_S
    =\bigotimes_{j\in S}q_j
     \otimes
     \bigotimes_{j\notin S}\gamma_j.
\]
Then
\begin{equation}
    D_{\mathrm{KL}}(\pi_C\|\pi_S)
    \ge (N-|S|)\kappa,
\end{equation}
and hence
\begin{equation}
    M_{\mathrm{post}}(\delta)
    \ge N-\delta/\kappa.
    \label{eq:T4-Mpost}
\end{equation}
If the BO utility depends only on a fixed set of blocks and its relevant
conditioning lies in that fixed set, then $M_{\mathrm{dec}}(\epsilon)=O(1)$
while $M_{\mathrm{post}}(\delta)=\Omega(N)$.

\begin{proof}
KL divergence is additive for product measures.  Therefore
\begin{align*}
D_{\mathrm{KL}}(\pi_C\|\pi_S)
&=
\sum_{j\notin S}D_{\mathrm{KL}}(q_j\|\gamma_j)\\
&\ge(N-|S|)\kappa.
\end{align*}
Rearranging the condition
$D_{\mathrm{KL}}(\pi_C\|\pi_S)\le\delta$ proves
Equation~\eqref{eq:T4-Mpost}.  The decision statement follows immediately when
the decision and its relevant factors are confined to a fixed block set.
\end{proof}

\paragraph{Interpretation discipline.}
T4A is the substantive result: it says that under a localized decision regime and
decaying structural influence, adding remote conditioning factors need not make
the BO decision harder.  T4B is only a clean witness that posterior-faithful
original-factor selection can behave very differently.  Do not market T4B as a
universal information-theoretic lower bound.

\paragraph{Main-paper use.}
Include a compact statement/diagram and proof intuition.  Put graph-tail details,
comparison-matrix inverse-decay conditions, and the product witness in the
supplement if page pressure is high.

% -----------------------------------------------------------------------------
\subsection{Inference contract and rigorous instantiation}

\paragraph{Role in the paper.}
Inference is a modular component of the certification framework, not a separate
methodological contribution.  T3 only requires a valid bound of the form
\[
    |\widehat G_S-G_S|\le B_{\mathrm{infer}}
\]
at a stated confidence level.

\paragraph{What is required for a rigorous experiment.}
The confidence statement must survive:
\begin{itemize}
    \item selection of the current leader;
    \item global challenger search;
    \item multiple adaptive factor-refinement rounds.
\end{itemize}
A safe finite-action construction is either (i) simultaneous acquisition/gap
confidence intervals with a union/confidence-sequence argument, or (ii)
independent certification samples after candidate selection, with explicit
allocation of the total failure probability across rounds.

\paragraph{Recommended rigorous backend.}
Use independent samples from the GP reference with importance weights
$w_S(f)=\exp[-E_S(f)]$ after shifting the energy so that a known upper bound on
the weights is available.  Estimate the numerator and normalizer separately and
construct a ratio interval.  For EI, the numerator is unbounded but has known or
bounded second moment under the Gaussian reference; use a finite-variance robust
mean estimator (e.g. median-of-means/Catoni) rather than an unjustified bounded
Hoeffding interval.  A lower confidence bound on the normalizer is required.  If
that lower bound is nonpositive, the certificate is inconclusive and inference
must be refined.

\paragraph{BLOCKER I1.}
Implement and validate one complete finite-sample inference-error construction
for a finite action grid.  Practical ESS, HMC/NUTS standard errors, SMC particle
diagnostics, or FlowGP/LatentFlow discretization diagnostics can remain
empirical elsewhere, but they do not substitute for this flagship certified
instantiation.

% -----------------------------------------------------------------------------
\subsection{Canonical algorithm implied by T1--T3}

\paragraph{Algorithmic invariant.}
At every refinement round, maintain an upper bound on the full-conditioned
regret of the current action.

\begin{enumerate}
    \item Infer under the active target $\pi_S$ and construct the required
          $B_{\mathrm{infer}}$.
    \item Optimize the active acquisition to obtain a leader $\widehat x$.
    \item Aggregate omitted factor sensitivities and construct
          $B_{\mathrm{struct}}(x,\widehat x)$.
    \item Solve one global optimistic challenger problem
          \[
          \max_x \Psi_S(x;\widehat x).
          \]
    \item If the certified upper bound is at most $\epsilon$, return
          $\widehat x$.
    \item Otherwise diagnose whether structural, inference, or challenger
          optimization error dominates.
    \item If structural error dominates, rank omitted factors by their
          contribution to the worst challenger's structural bound and activate
          a batch.  If inference or optimization dominates, refine that component
          without adding factors.
\end{enumerate}

The algorithm does not require all pairwise action comparisons.  It compares all
challengers to one current leader through a single optimistic envelope.

% -----------------------------------------------------------------------------
\subsection{Proof and presentation checklist}

\begin{description}
    \item[T1] Proof complete.  Main text.
    \item[T2 abstract] Proof complete.  Main text.
    \item[T2 reflection symmetry] Analytic construction complete under the
          stated assumptions.  \textbf{PROVED.}
    \item[T2 nonlinear PDE] Verify path-uniform constants and $A\succ0$ for the
          nonlinear-PDE factor family.  \textbf{BLOCKER.}
    \item[T3] Proof complete.  Main text.
    \item[T4] Abstract proof complete.  Clean up constants and connect the
          influence-decay assumptions to the actual factor graphs.
          \textbf{PROOF TASK.}
    \item[Inference] Produce one finite-sample, adaptivity-safe
          $B_{\mathrm{infer}}$.  \textbf{BLOCKER for an end-to-end certified claim.}
    \item[Continuous factors] Derive as a corollary only if used experimentally.
          Supplement.
\end{description}
